\documentclass[12pt,a4paper]{ctexart}
\usepackage[a4paper,margin=2cm]{geometry}
\usepackage{amsmath,amssymb,booktabs}

\begin{document}

\begin{center}
    {\LARGE \textbf{误差理论与数据处理}}
\end{center}

\section{基本概念}

\begin{equation*}
\text{绝对误差}=\text{测得值}-\text{真值}
\end{equation*}

\begin{equation*}
\text{修正值}=\text{真值}-\text{测得值}
\end{equation*}

\begin{equation*}
\text{相对误差}=\frac{\text{绝对误差}}{\text{真值}} \approx \frac{\text{绝对误差}}{\text{测得值}}
\end{equation*}

\begin{equation*}
\text{引用误差}=\frac{\text{绝对误差}}{\text{特定值}}=\frac{\text{示值误差}}{\text{测量范围上限}}
\end{equation*}

数字舍入规则：四舍六入五成双（舍入部分的数值而非下一位的数值）。

数据运算规则：
\begin{itemize}
\item 加减运算：以小数位数最少的数据位数为准，其他多取一位小数，最后结果与小数位数最少的数据小数位相同。
\item 乘除运算：以有效位数最少的数据位数为准，其他多取一位小数，最后结果与有效位数最少的数据小数位相同。
\item 平方开方：平方相当于乘法运算，开方等价于平方的逆运算。
\end{itemize}

\section{误差的基本性质与处理}

\subsection{随机误差}

残余误差校核：

\begin{enumerate}
    \item 残余误差代数和应符合：
    \begin{equation*}
    \begin{alignedat}{3}
        &\sum_{i=1}^{n} l_i = n\bar{x},\quad
        &&\bar{x}\text{ 为非凑整的准确数时，}\quad
        &&\sum_{i=1}^{n} v_i = 0;\\
        &\sum_{i=1}^{n} l_i > n\bar{x},\quad
        &&\bar{x}\text{ 为凑整的非准确数时，}\quad
        &&\sum_{i=1}^{n} v_i > 0,\ \text{其大小为求 }\bar{x}\text{ 时的余数};\\
        &\sum_{i=1}^{n} l_i < n\bar{x},\quad
        &&\bar{x}\text{ 为凑整的非准确数时，}\quad
        &&\sum_{i=1}^{n} v_i < 0,\ \text{其大小为求 }\bar{x}\text{ 时的亏数}.
    \end{alignedat}
    \end{equation*}

    \item 残余误差代数和绝对值应符合：
    \begin{equation*}
    \begin{alignedat}{2}
        &n\text{ 为偶数时，}\quad
        &&\left|\sum_{i=1}^{n} v_i\right| \leq \frac{n}{2}A;\\
        &n\text{ 为奇数时，}\quad
        &&\left|\sum_{i=1}^{n} v_i\right| \leq \left(\frac{n}{2}-0.5\right)A.
    \end{alignedat}
    \end{equation*}
    式中的 $A$ 为实际求得的算术平均值 $\bar{x}$ 末位数的一个单位。
\end{enumerate}

等精度测量中，单次测量的标准差如下，其中 $n$ 为测量次数，$\delta_i$ 为测量值与真值之差。
\begin{equation*}
\sigma=\sqrt{\frac{\sum_{i=1}^{n}\delta_i^2}{n}}
\end{equation*}

被测量的真值未知时，采用贝塞尔公式，其中 $v_i$ 为残余误差。
\begin{equation*}
\sigma=\sqrt{\frac{\sum_{i=1}^{n}v_i^2}{n-1}}
\end{equation*}

或然误差
\begin{equation*}
\rho=\frac{2}{3}\sqrt{\frac{\sum_{i=1}^{n}v_i^2}{n-1}}
\end{equation*}

平均误差
\begin{equation*}
\theta=\frac{4}{5}\sqrt{\frac{\sum_{i=1}^{n}v_i^2}{n-1}}
\end{equation*}

\subsection{算术平均值误差评定}

算术平均值的标准差
\begin{equation*}
\sigma_{\bar{x}}=\frac{\sigma}{\sqrt{n}}
\end{equation*}

或然误差 $R$
\begin{equation*}
R=\frac{\rho}{\sqrt{n}}=\frac{2}{3}\sqrt{\frac{\sum_{i=1}^{n}v_i^2}{n(n-1)}}
\end{equation*}

平均误差 $T$
\begin{equation*}
T=\frac{\theta}{\sqrt{n}}=\frac{4}{5}\sqrt{\frac{\sum_{i=1}^{n}v_i^2}{n(n-1)}}
\end{equation*}

标准差计算：别捷尔斯法
\begin{equation*}
\sigma=1.253\,\frac{\sum_{i=1}^{n}\lvert v_i\rvert}{\sqrt{n(n-1)}}
\end{equation*}

标准差计算：极差法（$n<10$）
\begin{equation*}
\sigma=\frac{x_{\max}-x_{\min}}{d_n}
\end{equation*}

标准差计算：最大误差法
\begin{equation*}
\sigma=\dfrac{{\lvert \delta_i\rvert}_{\max}}{K_n}=\frac{{\lvert v_i\rvert}_{\max}}{K_n^{\prime}}
\end{equation*}

\subsection{不等精度测量}

$t$分布计算测量列算术平均值的极限误差：
\begin{equation*}
\delta_{\lim}\bar{x}=\pm t_{\alpha} \sigma_{\bar{x}}
\end{equation*}

权的计算：
\begin{equation*}
p_1:p_2:\ldots:p_m=\frac{1}{\sigma_{\bar{x}_1}^2}:\frac{1}{\sigma_{\bar{x}_2}^2}:\ldots:\frac{1}{\sigma_{\bar{x}_m}^2}
\end{equation*}

加权算术平均值：
\begin{equation*}
    \bar{x}=x_0+\frac{\displaystyle\sum_{i=1}^{m}p_i(x_i-x_0)}
    {\displaystyle\sum_{i=1}^{m}p_i}
\end{equation*}

不等精度测量时的算术平均值的标准差（$\sigma$已知）：
\begin{equation*}
\sigma_{\bar{x}}=\frac{\sigma}{\sqrt{\sum_{i=1}^m p_i}}
\end{equation*}

不等精度测量时的算术平均值的标准差（$\sigma$未知）：
\begin{equation*}
\sigma_{\bar{x}}=\sqrt{\dfrac{\sum_{i=1}^{m}p_i\,v_{\bar{x}_i}^2}{(m-1)\sum_{i=1}^{m}p_i}}
\end{equation*}

\begin{table}[htbp]
\centering
\caption{常见误差分布的分布密度与标准差}
\renewcommand{\arraystretch}{2.0}
\setlength{\tabcolsep}{16pt}
\begin{tabular}{lll}
\hline
分布 & 分布密度 $f(\delta)$ & 标准差 $\sigma$ \\
\hline
正态
&
$\displaystyle f(\delta)=\frac{1}{\sigma\sqrt{2\pi}}e^{-\frac{\delta^2}{2\sigma^2}}$
&
$\displaystyle \sigma$
\\
均匀
&
$\displaystyle f(\delta)=\frac{1}{2a}\quad |\delta|\le a$
&
$\displaystyle a/\sqrt{3}$
\\
反正弦
&
$\displaystyle f(\delta)=\frac{1}{\pi\sqrt{a^2-\delta^2}}\quad |\delta|<a$
&
$\displaystyle a/\sqrt{2}$
\\
三角
&
$\displaystyle
f(\delta)=
\begin{cases}
\dfrac{a+\delta}{a^2}, & -a\le \delta<0 \\[8pt]
\dfrac{a-\delta}{a^2}, & 0\le \delta\le a
\end{cases}$
&
$\displaystyle a/\sqrt{6}$
\\
直角分布
&
$\displaystyle f(\delta)=\frac{\delta+a}{2a^2}\quad |\delta|\le a$
&
$\displaystyle \sqrt{2}\,a/3$
\\
椭圆分布
&
$\displaystyle f(\delta)=\frac{2}{\pi a^2}\sqrt{a^2-\delta^2}\quad |\delta|\le a$
&
$\displaystyle a/2$
\\
双三角
&
$\displaystyle f(\delta)=\frac{|\delta|}{a^2}\quad |\delta|\le a$
&
$\displaystyle a/\sqrt{2}$
\\
\hline
\end{tabular}
\end{table}

\newpage

\subsection{系统误差的判别}

若满足以下公式，则认为存在系统误差。

不同公式计算标准差比较法：
\begin{equation*}
\left|\frac{\sigma_2}{\sigma_1}-1\right|\geq \frac{2}{\sqrt{n-1}}
\end{equation*}

计算数据比较法：
\begin{equation*}
|\bar{x}_i-\bar{x}_j|<2\sqrt{\sigma_i^2+\sigma_j^2}
\end{equation*}

秩和检验法，$n_1,n_2\leq 10$ 时（同则平均）：
\begin{equation*}
T\notin(T_-,T_+)
\end{equation*}

秩和检验法，$n_1,n_2>10$ 时：
\begin{equation*}
\left|\frac{T-a}{\sigma}\right|
=
\left|
\frac{T-\dfrac{n_1(n_1+n_2+1)}{2}}
{\sqrt{\dfrac{n_1n_2(n_1+n_2+1)}{12}}}
\right|
>t_{\alpha}
\end{equation*}

$t$ 检验法
\begin{equation*}
|t|
=
(\bar{x}-\bar{y})\,
\sqrt{
\frac{n_x n_y (n_x+n_y-2)}
{(n_x+n_y)(n_x \sigma_x^2+n_y \sigma_y^2)}
}
\geq t_\alpha
\end{equation*}

\subsection{粗大误差的判别}

若满足以下公式，则认为存在粗大误差。

莱以特准则（$3\sigma$ 准则）：
\begin{equation*}
|v_i|>3\sigma=3\sqrt{\frac{\sum_{i=1}^{n}v_i^2}{n-1}}
\end{equation*}

罗曼诺夫斯基准则：
\begin{equation*}
\frac{|x_j-\bar{x}|}{\sigma}
=
\frac{
\left|x_j-\dfrac{1}{n-1}\sum_{i=1,i\neq j}^{n}x_i\right|
}{
\sqrt{\dfrac{\sum_{i=1}^{n}v_i^2}{n-2}}
}
>K(n,\alpha)
\end{equation*}

格罗布斯准则：
\begin{equation*}
g_{(1)}=\frac{\bar{x}-x_{(1)}}{\sigma}
\quad
g_{(n)}=\frac{x_{(n)}-\bar{x}}{\sigma}
\quad
g_{(i)}\geq g_0(n,\alpha)
\end{equation*}

狄克松准则（不需要计算标准差）：
\begin{equation*}
\begin{cases}
r_{10}=\dfrac{x_{(2)}-x_{(1)}}{x_{(n)}-x_{(1)}} =\dfrac{x_{(n)}-x_{(n-1)}}{x_{(n)}-x_{(1)}} & > r_0(n,\alpha),\quad n\leq 7 \\[8pt]
r_{11}=\dfrac{x_{(2)}-x_{(1)}}{x_{(n-1)}-x_{(1)}} =\dfrac{x_{(n)}-x_{(n-1)}}{x_{(n)}-x_{(2)}} & > r_0(n,\alpha),\quad 8\leq n\leq 10 \\[8pt]
r_{21}=\dfrac{x_{(3)}-x_{(1)}}{x_{(n-1)}-x_{(1)}} =\dfrac{x_{(n)}-x_{(n-2)}}{x_{(n)}-x_{(2)}} & > r_0(n,\alpha),\quad 11\leq n\leq 13 \\[8pt]
r_{22}=\dfrac{x_{(3)}-x_{(1)}}{x_{(n-2)}-x_{(1)}} =\dfrac{x_{(n)}-x_{(n-2)}}{x_{(n)}-x_{(3)}} & > r_0(n,\alpha),\quad n\geq 14
\end{cases}
\end{equation*}

\section{误差合成与分配}

函数的系统误差
\begin{equation*}
\Delta y = \frac{\partial f}{\partial x_1} \Delta x_1 + \frac{\partial f}{\partial x_2} \Delta x_2 + \cdots + \frac{\partial f}{\partial x_n} \Delta x_n
\end{equation*}

函数的标准差合成
{\small
\begin{equation*}
\sigma_y=
\sqrt{
\left(\frac{\partial f}{\partial x_1}\right)^2 \sigma_{x_1}^2
+
\left(\frac{\partial f}{\partial x_2}\right)^2 \sigma_{x_2}^2
+ \cdots +
\left(\frac{\partial f}{\partial x_n}\right)^2 \sigma_{x_n}^2
+
2 \sum_{1 \le i < j}^{n}
\left(
\frac{\partial f}{\partial x_i}
\frac{\partial f}{\partial x_j}
\rho_{ij}\sigma_{x_i}\sigma_{x_j}
\right)
}
\end{equation*}
}

函数的极限误差
\begin{equation*}
\delta_{\lim y} = \pm \sqrt{\delta_{\lim x_1}^{2}+\delta_{\lim x_2}^{2}+\cdots+\delta_{\lim x_n}^{2}}
\end{equation*}

按等作用原则进行误差分配，各分量标准差
\begin{equation*}
\sigma_i=\frac{\sigma_y}{\sqrt{n}}\frac{1}{\partial f/\partial x_i}
=\frac{\sigma_y}{\sqrt{n}}\frac{1}{a_i}
\end{equation*}

按等作用原则进行误差分配，各分量极限误差
\begin{equation*}
\delta_i=\frac{\delta}{\sqrt{n}}\frac{1}{\partial f/\partial x_i}
=\frac{\delta}{\sqrt{n}}\frac{1}{a_i}
\end{equation*}

\section{不确定度}

根据分布区间半宽$a$和包含因子$k_p$估计标准不确定度

\begin{table}[htbp]
\centering
\caption{标准不确定度的常用估计公式}
\begin{tabular}{ll}
\toprule
分布类型/已知条件 & 标准不确定度 \\
\midrule
正态分布 & $u_x = a/k_p$ \\
已知扩展不确定度 & $u_x = U_x/k$ \\
均匀分布 & $u_x = a/\sqrt{3}$ \\
三角分布 & $u_x = a/\sqrt{6}$ \\
反正弦分布 & $u_x = a/\sqrt{2}$ \\
\bottomrule
\end{tabular}
\end{table}

A类评定时，自由度（最大误差法，$n=1$）
\begin{equation*}
\nu=0.9
\end{equation*}

A类评定时，自由度（贝塞尔公式）
\begin{equation*}
\nu=n-1
\end{equation*}

B类评定时，自由度
\begin{equation*}
\nu=\frac{1}{2\left(\sigma_u/u\right)^2}
\end{equation*}

不确定度的合成
\begin{equation*}
u_c=\sqrt{\sum_{i=1}^{N}\left(\frac{\partial f}{\partial x_i}\right)^2u_{x_i}^2}
=\sqrt{\sum_{i=1}^{N}u_i^2}
\end{equation*}

展伸不确定度
\begin{equation*}
U=ku_c
\end{equation*}

自由度合成
\begin{equation*}
\nu=\frac{U_c^{4}}{\sum_{i=1}^{n} \frac{u_i^4}{v_i}}
\end{equation*}

\section{最小二乘法}

线性测量参数的误差方程式
\begin{equation*}
V=L-A\hat{X}
\end{equation*}

正规方程的解（等精度）
\begin{equation*}
X=C^{-1} A^T L={(A^T A)}^{-1} A^T L
\end{equation*}

正规方程的解（不等精度）
\begin{equation*}
X=C^{-1} A^T L={(A^T PA)}^{-1} A^T PL
\end{equation*}

标准差（$n$ 为方程个数，$t$ 为未知量个数）（等精度）
\begin{equation*}
\sigma=\sqrt{\frac{\sum_{i=1}^{n}v_i^2}{n-t}}
\end{equation*}

标准差（$n$ 为方程个数，$t$ 为未知量个数）（不等精度）
\begin{equation*}
\sigma=\sqrt{\frac{\sum_{i=1}^{n}p_i v_i^2}{n-t}}
\end{equation*}

各估计量的标准差
\begin{equation*}
\left.
\begin{aligned}
\sigma_{x_1} &= \sigma \sqrt{d_{11}} \\
\sigma_{x_2} &= \sigma \sqrt{d_{22}} \\
&\ \ \vdots \\
\sigma_{x_t} &= \sigma \sqrt{d_{tt}}
\end{aligned}
\right\}
\end{equation*}

协方差矩阵
\begin{equation*}
C^{-1}
=
\begin{pmatrix}
d_{11} & d_{12} & \cdots & d_{1t} \\
d_{21} & d_{22} & \cdots & d_{2t} \\
\vdots & \vdots & \ddots & \vdots \\
d_{t1} & d_{t2} & \cdots & d_{tt}
\end{pmatrix}
\end{equation*}

\section{回归分析}

设一元线性回归方程为
\begin{equation*}
\hat{y}=b_0+bx
\end{equation*}

其中，$b$ 为回归系数，$b_0$ 为截距。

样本均值
\begin{equation*}
\bar{x}=\frac{1}{N}\sum_{t=1}^{N}x_t
\end{equation*}

\begin{equation*}
\bar{y}=\frac{1}{N}\sum_{t=1}^{N}y_t
\end{equation*}

回归系数
\begin{equation*}
b=\frac{l_{xy}}{l_{xx}}
\end{equation*}

截距
\begin{equation*}
b_0=\bar{y}-b\bar{x}
\end{equation*}

自变量离差平方和
\begin{equation*}
l_{xx}=\sum_{t=1}^{N}(x_t-\bar{x})^2
=\sum_{t=1}^{N}x_t^2-\frac{1}{N}\left(\sum_{t=1}^{N}x_t\right)^2
\end{equation*}

离差积和
\begin{equation*}
l_{xy}=\sum_{t=1}^{N}(x_t-\bar{x})(y_t-\bar{y})
=\sum_{t=1}^{N}x_t y_t-\frac{1}{N}\left(\sum_{t=1}^{N}x_t\right)\left(\sum_{t=1}^{N}y_t\right)
\end{equation*}

因变量离差平方和
\begin{equation*}
l_{yy}=\sum_{t=1}^{N}(y_t-\bar{y})^2
=\sum_{t=1}^{N}y_t^2-\frac{1}{N}\left(\sum_{t=1}^{N}y_t\right)^2
\end{equation*}

总离差平方和
\begin{equation*}
S=l_{yy}
\end{equation*}

回归平方和
\begin{equation*}
U=b l_{xy}
\end{equation*}

残差平方和
\begin{equation*}
Q=l_{yy}-bl_{xy}
\end{equation*}

平方和关系
\begin{equation*}
S=U+Q
\end{equation*}

统计量
\begin{equation*}
F=\frac{U/\nu_U}{Q/\nu_Q}=\frac{U/1}{Q/(N-2)}
\end{equation*}

自由度
\begin{equation*}
\nu_S=N-1, \qquad \nu_U=1, \qquad \nu_Q=N-2
\end{equation*}

残余标准差
\begin{equation*}
\sigma=\sqrt{\frac{Q}{N-2}}
\end{equation*}

相关指数
\begin{equation*}
R^2=1-\frac{\sum_{t=1}^{N}(y_t-\hat{y}_t)^2}{\sum_{t=1}^{N}(y_t-\bar{y})^2}
=1-\frac{Q}{S}
\end{equation*}

\begin{table}[h]
\centering
\renewcommand{\arraystretch}{1.3}
\caption{回归方差分析表}
\begin{tabular}{cccccc}
\toprule
来源 & 平方和 & 自由度 & 方差 & $F$ & 显著性 \\
\midrule
回归 &
$U=b\,l_{xy}$ &
$1$ &
-- &
$\dfrac{U/1}{Q/(N-2)}$ &
$\alpha$ \\

残余 &
$Q=l_{yy}-b\,l_{xy}$ &
$N-2$ &
$Q/(N-2)$ &
-- &
-- \\
\midrule
总计 &
$S=l_{yy}$ &
$N-1$ &
-- &
-- &
-- \\
\bottomrule
\end{tabular}
\end{table}

\end{document}